% Bounded development: interval energy reader of the register K.
% Status: added formalization on K and the Schur identities
% already present in the corpus. No historical novelty is attributed to Schur.
\section{Dirichlet forms and variational equivalence of refinement}
\label{sec:energia-refinamiento-k}

The positive representation of the dodecaphasic register defines an
ordered partition through
\[
 K=(234,543,140,729,659,824,621,58,914,794,146,601),
 \qquad d_j=K_j/6263,
 \qquad t_i=\sum_{j=1}^id_j.
\]
The register is received from the composition
\(x_{\rm term}\mapsto\mathcal R_{12}\mapsto
(b^{90},b^{120},Q_{\rm TPK})\mapsto U_{\rm sgn}\mapsto K\),
subsequent to APP--TRIT--TPK and its enriched state, in accordance
with the readers reconstructed in \cite{hmtX}. This section applies
a quadratic reader to the resulting partition. Depths will be compared
through restrictions to inherited endpoints, so that every energy
coefficient retains its origin in an already produced positive separation.

The corpus contains a weighted difference energy for the graph of
the two memories and a Schur reduction of the complete forms.
Here these operations are brought together on the interval network
of the register. The specialization to conductances \(1/d_j\),
the interpolation operator, and its compatibility with deformation
of the separations constitute the formalization developed in this section.

\subsection{Reconstruction of separations from the stiffness matrix}

Let \(d=(d_1,\ldots,d_m)\), with \(d_j>0\), and let
\(f=(f_0,\ldots,f_m)\in\mathbb C^{m+1}\) be the values of an
observable at the partition endpoints. Define
\begin{equation}
 \mathcal E_d(f)=\sum_{j=1}^{m}\frac{|f_j-f_{j-1}|^2}{d_j},
 \qquad L_d=D^*\operatorname{diag}(d_1^{-1},\ldots,d_m^{-1})D,
 \label{eq:energia-intervalar}
\end{equation}
where \(D:\mathbb C^{m+1}\to\mathbb C^m\) is the oriented incidence
operator, \((Df)_j=f_j-f_{j-1}\), and the asterisk denotes the
Hermitian adjoint. Thus \(\mathcal E_d(f)=f^*L_df\).
The stiffness matrix has entries
\[
 (L_d)_{00}=d_1^{-1},\quad (L_d)_{mm}=d_m^{-1},\quad
 (L_d)_{ii}=d_i^{-1}+d_{i+1}^{-1}\quad(1\le i<m),
\]
\[
 (L_d)_{i-1,i}=(L_d)_{i,i-1}=-d_i^{-1},
\]
and all other entries vanish. It is Hermitian positive semidefinite,
its kernel consists of the constant functions, and its form is positive
definite on the quotient by that kernel. The complete matrix recovers
the separations through
\begin{equation}
 d_j=-\frac1{(L_d)_{j-1,j}}.
 \label{eq:recuperacion-d-rigidez}
\end{equation}

\begin{proposition}[Equivalence between separations and the complete nodal reader]
\label{prop:energia-bidireccional}
The map \(d\mapsto L_d\) is a bijection between
\((0,\infty)^m\) and the set of real symmetric tridiagonal matrices
whose entries immediately adjacent to the diagonal are negative and
whose row sums are zero. Its inverse is
\eqref{eq:recuperacion-d-rigidez}. The normalization
\(\sum_jd_j=1\) is equivalent to
\(\sum_{j=1}^m-1/(L_d)_{j-1,j}=1\).
\end{proposition}
\begin{proof}
The explicit formulas for \(L_d\) satisfy all the conditions.
Conversely, the negative adjacent entries determine unique positive
separations through \eqref{eq:recuperacion-d-rigidez}. The zero
sum of each row determines its diagonal entries as the sum of
incident conductances. Tridiagonality makes all remaining entries
zero. The matrix in \eqref{eq:energia-intervalar} is recovered
exactly; its positive semidefiniteness then follows from the
factorization \(D^*\operatorname{diag}(1/d_j)D\).
\end{proof}

The weighted discrete derivative \(J_j=(f_j-f_{j-1})/d_j\) is the
oriented flow of edge \(j\). At interior vertices,
\((L_df)_i=J_i-J_{i+1}\); at the endpoints,
\((L_df)_0=-J_1\) and \((L_df)_m=J_m\). The interior equation
\(L_df=0\) expresses that the flow has the same value on every edge.
This defines a quadratic observable of the interval reader;
its realization as a physical quantity requires the relevant dimensional chart.

\subsection{Harmonic interpolation, Schur complement, and detail reconstruction}

Refine each separation as
\(d_j=\sum_{c=1}^{r_j}\delta_{j,c}\), with \(\delta_{j,c}>0\).
Let \(\widetilde d\) be the concatenation of these separations,
\(R\) restriction to the old endpoints, and \(H\) the interpolation
defined by
\begin{equation}
 (Hf)_{j,c}=(1-\theta_{j,c})f_{j-1}+\theta_{j,c}f_j,
 \qquad
 \theta_{j,c}=\frac{\sum_{a=1}^c\delta_{j,a}}{d_j}.
 \label{eq:interpolacion-armonica-k}
\end{equation}
Shared endpoints are counted only once. We have \(RH=I\).

\begin{theorem}[Isometric refinement of the energy]
\label{thm:energia-schur-refinamiento}
For every vector of old values \(f\),
\begin{equation}
 H^*L_{\widetilde d}H=L_d,
 \qquad L_{\widetilde d}H=R^*L_d,
 \qquad
 \min_{Rg=f}\mathcal E_{\widetilde d}(g)=\mathcal E_d(f).
 \label{eq:isometria-energia-k}
\end{equation}
The minimizer is unique and equals \(Hf\). Every \(g\) admits
the unique decomposition
\begin{equation}
 g=H(Rg)+z,\qquad Rz=0,
 \qquad
 \mathcal E_{\widetilde d}(g)
 =\mathcal E_d(Rg)+\mathcal E_{\widetilde d}(z).
 \label{eq:detalle-energia-k}
\end{equation}
\end{theorem}
\begin{proof}
On an old edge, write \(h_c=g_{j,c}-g_{j,c-1}\), so that
\(\sum_ch_c=f_j-f_{j-1}=a\). If
\(h_c=\delta_{j,c}a/d_j+r_c\), then \(\sum_cr_c=0\) and
\[
 \sum_c\frac{|h_c|^2}{\delta_{j,c}}
 =\frac{|a|^2}{d_j}+\sum_c\frac{|r_c|^2}{\delta_{j,c}}.
\]
The equality follows by expansion; the cross term is proportional to
\(\sum_cr_c\) and vanishes. The minimum is attained exactly when
all \(r_c\) vanish, the condition determining
\eqref{eq:interpolacion-armonica-k}. Summing over the edges gives
the minimum and the energy decomposition. The flow of \(Hf\) is
constant within every old edge, so its interior components of
\(L_{\widetilde d}Hf\) vanish and its inherited components agree
with \(L_df\). This proves the second matrix identity; multiplying
by \(H^*\), using \(RH=I\), gives the first.
\end{proof}

The term \(\mathcal E_{\widetilde d}(z)\) is nonnegative and
vanishes exactly for \(z=0\): a zero-energy detail would be constant
and, vanishing at the inherited endpoints, would vanish throughout
the network. The pair \((Rg,z)\) preserves the complete refined
vector and allows its recovery through \(g=H(Rg)+z\).
Minimization selects its harmonic component; storing the detail \(z\)
retains the information that the reduced form alone does not distinguish.
Conservation of the observed energy and conservation of the refined
state are thus expressed by different explicit maps.

Ordering vertices as inherited endpoints \(E\) and interior vertices
\(I\), write
\[
 L_{\widetilde d}=\begin{pmatrix}A&B^*\\B&C\end{pmatrix}.
\]
For a refinement with interior vertices, the block \(C\) is
positive definite: a function vanishing at the endpoints and having
zero energy is constant on every subinterval and therefore identically
zero. We obtain
\begin{equation}
 H=\begin{pmatrix}I\\-C^{-1}B\end{pmatrix},\qquad
 \operatorname{Schur}_{I}(L_{\widetilde d})
 =A-B^*C^{-1}B=L_d.
 \label{eq:schur-intervalar-k}
\end{equation}
If refinement preserves all vertices and introduces none, the same
statement reduces to \(H=I\) and \(L_{\widetilde d}=L_d\).

\begin{proposition}[Naturality of prolongations and elimination]
\label{prop:energia-naturalidad-sucesiva}
For three nested partitions, harmonic interpolations and Schur
complements satisfy
\begin{equation}
 H_{n+2\leftarrow n}
 =H_{n+2\leftarrow n+1}H_{n+1\leftarrow n},
 \qquad
 \operatorname{Schur}_{n+1\to n}
 \bigl(\operatorname{Schur}_{n+2\to n+1}(L_{n+2})\bigr)=L_n.
 \label{eq:energia-naturalidad-sucesiva}
\end{equation}
\end{proposition}
\begin{proof}
Interpolation of affine values on an interval remains the same affine
function after any subdivision of that interval. This proves the
first identity. Minimizing first over vertices added at the last
level and then over those at the intermediate level amounts to
minimizing over all of them, with the same inherited endpoints fixed.
The variational identity in \eqref{eq:isometria-energia-k} and
uniqueness of the minimizer prove the second identity. Restriction
of the minimum is therefore performed through the same composition
as harmonic prolongation of the data.
\end{proof}

\subsection{Dirichlet-to-Neumann response and energy defect of an extension}

If only the initial and final endpoints are preserved, let
\(D_0=\sum_{j=1}^md_j\). For values \((a,b)\), the harmonic
solution is
\[
 f_i=a+(b-a)\frac{\sum_{j=1}^id_j}{D_0},\qquad
 J_j=(b-a)/D_0.
\]
Its Dirichlet-to-Neumann operator, taking boundary values to outgoing
flows, is
\begin{equation}
 \Lambda_d=\frac1{D_0}
 \begin{pmatrix}1&-1\\-1&1\end{pmatrix},\qquad
 \min\mathcal E_d=\frac{|b-a|^2}{D_0}.
 \label{eq:dtn-intervalar-k}
\end{equation}
When values are prescribed at all inherited vertices, the same
argument gives the multipoint response operator \(\Lambda=L_d\).

An approximate extension \(g=(f,y)\), where \(y\) collects the
values at interior vertices, has interior residual \(r=Bf+Cy\).
The exact identity
\begin{equation}
 \mathcal E_{\widetilde d}(g)-f^*\Lambda f
 =r^*C^{-1}r\ge0
 \label{eq:residual-dtn-k}
\end{equation}
follows from \(y+C^{-1}Bf=C^{-1}r\) and the decomposition
\[
 g^*L_{\widetilde d}g
 =f^*(A-B^*C^{-1}B)f
 +(y+C^{-1}Bf)^*C(y+C^{-1}Bf).
\]
Thus the effective value and the detail retain a complete Hermitian
decomposition. In particular, if a linear extension is
\(\widetilde H=\begin{pmatrix}I\\Z\end{pmatrix}\), its residual
matrix is \(\mathscr R=B+CZ\), and
\[
 \widetilde H^*L_{\widetilde d}\widetilde H-\Lambda
 =\mathscr R^*C^{-1}\mathscr R\succeq0.
\]
A zero residual characterizes exact harmonic interpolation.
This energy residual lives in the space of interior vertices;
cancellation of residues of meromorphic forms belongs to another
boundary map and retains its own definition.

\subsection{Commutation of refinement and the separation flow}

Let \(u=P_3P_{11}K\) be the algebraic direction already recovered
from the register readers \cite{hmtX}. For the real parameter \(s\),
define
\[
 d_j(s)=\frac{K_je^{su_j}}{\sum_\ell K_\ell e^{su_\ell}}.
\]
Fix positive fractions \(\vartheta_{j,c}\) summing to one for each
\(j\). The inherited refinement is
\[
 \delta_{j,c}(s)=\vartheta_{j,c}d_j(s),\qquad
 u_{j,c}=u_j.
\]
First producing the refined weights \(K_j\vartheta_{j,c}\) and
then deforming them gives the same collection: the successors have identical
exponential factors, and their coefficients sum to \(K_j\).
For every \(s\),
\begin{equation}
 \operatorname{Schur}_{I}(L_{\widetilde d(s)})=L_{d(s)},
 \qquad H^*L_{\widetilde d(s)}H=L_{d(s)}.
 \label{eq:schur-flujo-k}
\end{equation}
In this parent--children comparison, \(H\) is constant: its
coefficients are the partial sums of \(\vartheta_{j,c}\).
By contrast, interpolation from the two global endpoints uses
\(t_i(s)=\sum_{j\le i}d_j(s)\) and changes with the parameter.
Both statements correspond to different restrictions of the same network.

The normalization \(\sum_jd_j(s)=1\) makes the global two-endpoint
operator remain equal to
\(\left(\begin{smallmatrix}1&-1\\-1&1\end{smallmatrix}\right)\).
The all-node matrix \(L_{d(s)}\) does preserve the deformation and
recovers it through \eqref{eq:recuperacion-d-rigidez}. The information
loss in the two-endpoint reader is identified exactly: it retains the
total length and its flows, whereas the nodal reader retains all
separations. The detail archive provides the complement needed to
reconstruct the refined observable.
